% \section{Proposed Research}

% We target a distributed client-server usage model where \textit{clients} run programs (e.g., desktops, servers, phones). The \textit{server} is where the in-house analysis that reproduces the bugs occurs. The server is a distributed system  (e.g., like the Windows Error Reporting System~\cite{wer}), but we refer to it as the \textit{server} for brevity. \textit{We target fail-stop bugs} (e.g., crash, hang, assertion failure, etc.), however developers can arbitrarily configure the clients to ship execution information to the server. Clients will have driver and runtime support to monitor executions as per the server's directions. The server will aggregate and process clients' data to reproduce bugs. The server will adapt the monitored execution information at the clients to improve the effectiveness and the accuracy of bug reproduction. Most operating systems have error reporting frameworks (Apport in Ubuntu~\cite{ubuntuApport}, WER in Windows~\cite{wer}, Crashlytics in iOS and Android~\cite{crashlytics}, CrashReporter in MacOSX~\cite{crashReporterMac}), which allows developers to receive reports (e.g., coredumps) when programs fail in production. We will target Apport~\cite{ubuntuApport} for the initial development and evaluation of our techniques and systems. We will then use Windows Error Reporting~\cite{wer} (where the current state-of-the-art failure analysis system, REPT, operates) for technology transfer and \textit{large-scale evaluation} (see letter from Dr. Cui at Microsoft).


